%% =========================================================
\section{Fixed-point algebra}
\label{sec:fixed_point_algebra}
%% =========================================================

\begin{definition}[Kraus fixed points]\label{def:kraus_fixed_points}
    \lean{Kraus.fixedPoints}\leanok
    \uses{def:kraus_map}
    The \emph{fixed-point set} of a Kraus map
    $E(X)=\sum_iK_iXK_i^\dagger$ is
    $\Fix(E)=\{X\in\MN{D}:E(X)=X\}$.
\end{definition}

\begin{definition}[Fixed points of a linear map]\label{def:abstract_fixed_points}
    \lean{SchwarzMap.fixedPoints}\leanok
    The fixed-point set of a complex-linear map
    $E:\MN{D}\to\MN{D}$ is $\Fix(E)=\{X\in\MN{D}:E(X)=X\}$.
\end{definition}

\begin{theorem}[Fixed points lie in the multiplicative domain]
    \label{thm:fixedPoints_in_multiplicativeDomain}
    \lean{Kraus.fixedPoints_in_multiplicativeDomain}
    \leanok
    \uses{def:kraus_fixed_points, def:kraus_adjoint_map, def:tp_kraus,
        def:full_multiplicative_domain}
    Let $E$ be a trace-preserving Kraus map whose Schr{\"o}dinger-picture map has a
    positive definite fixed point $\rho>0$.
    Then every adjoint fixed point belongs to the multiplicative domain of the
    adjoint Kraus family.
\end{theorem}

\begin{proof}\leanok
    The Kadison--Schwarz equality for fixed points forces both one-sided
    multiplicative identities, so the fixed point lies in the full
    multiplicative domain.
\end{proof}

\begin{theorem}[Fixed-point algebra after choosing a faithful invariant weight]
    \label{thm:fixedPointsStarSubalgebra}
    \lean{SchwarzMap.fixedPointsStarSubalgebra}\leanok
    \uses{def:abstract_fixed_points, def:abstract_schwarz_inequality,
        lem:posDef_weighted_trace_faithful}
    Let $E:\MN{D}\to\MN{D}$ be positive and unital, and suppose that, for
    every $A\in\MN{D}$, it satisfies the Schwarz inequality
    \begin{align}
        E(A^\dagger A)-E(A^\dagger)E(A)
         & \succeq0.
        \label{eq:channel_fixed_point_schwarz}
    \end{align}
    If the trace-pairing adjoint $E^*$ has a positive definite fixed point
    $\rho>0$, then $\Fix(E)$ is a $*$-subalgebra of $\MN{D}$.

    In the terminology introduced immediately before
    \cite[Example~5.3]{Wolf2012Quantum}, a \emph{Schwarz map} is precisely a
    positive unital map satisfying (\ref{eq:channel_fixed_point_schwarz}).
    Thus the three assumptions reproduce the source's Schwarz-map convention.
    In \cite[Theorem~6.12]{Wolf2012Quantum}, the trace adjoint is initially
    assumed to have an arbitrary full-rank fixed point. The proposition on
    positive fixed points then produces a positive definite fixed point
    $\rho$, to which the theorem above applies. The resulting explicit block
    form of the algebra is given in
    \cite[Equation~(1.39)]{Wolf2012Quantum}.
\end{theorem}

\begin{proof}\leanok
    \uses{thm:abstract_md_characterization}
    Positivity gives $E(X^\dagger)=E(X)^\dagger$, so fixed points are closed
    under adjoints. For a fixed point $X$, set
    $\Delta_X=E(X^\dagger X)-E(X^\dagger)E(X)\succeq0$.
    The trace-pairing identity, $E^*(\rho)=\rho$, and $E(X)=X$ give
    \begin{align}
        \tr(\rho\Delta_X)
         & =\tr\!\left(E^*(\rho)X^\dagger X\right)
        -\tr\!\left(\rho X^\dagger X\right)
        =0.
        \notag
    \end{align}
    Since $\rho>0$, Lemma~\ref{lem:posDef_weighted_trace_faithful}
    forces $\Delta_X=0$. Applying the same argument to $X^\dagger$ gives
    equality in the other one-sided Schwarz inequality. Hence
    Theorem~\ref{thm:abstract_md_characterization} places $X$ in the full
    multiplicative domain. Consequently, for fixed points $X,Y$,
    $E(XY)=E(X)E(Y)=XY$. Linearity and unitality supply the remaining
    subalgebra axioms.
\end{proof}

\begin{remark}[Kraus specialization]
    \lean{Kraus.fixedPointsStarSubalgebra}\leanok
    For a unital Kraus map whose adjoint has a positive definite fixed point,
    the weighted Kadison--Schwarz equality gives directly that its fixed
    points form a $*$-subalgebra.
\end{remark}

\begin{remark}\leanok
    Theorem~\ref{thm:fixedPoints_in_multiplicativeDomain} is the
    trace-preserving Heisenberg-picture counterpart of the same
    multiplicative-domain argument.
\end{remark}

\begin{theorem}[Projected elements lie in the multiplicative domain]
    \label{thm:map_mem_multiplicativeDomain_of_idempotent}
    \lean{Kraus.map_mem_multiplicativeDomain_of_idempotent}\leanok
    \uses{def:kraus_fixed_points, def:kraus_adjoint_map, def:unital_kraus,
        def:full_multiplicative_domain}
    Let $E$ be a unital Kraus map whose adjoint map has a positive definite
    fixed point, and suppose $E^2=E$. Then $E(X)$ lies in the multiplicative
    domain of $E$ for every $X\in\MN{D}$.
\end{theorem}

\begin{proof}\leanok
    \uses{thm:fixedPointsStarSubalgebra}
    Idempotence gives $E(E(X))=E(X)$, so $E(X)$ is a fixed point.
    By the argument in Theorem~\ref{thm:fixedPointsStarSubalgebra}, the
    Kadison--Schwarz gap for the fixed point $E(X)$ vanishes, hence
    $E(E(X)^\dagger E(X))=E(X)^\dagger E(X)$.
    Fixed points are closed under adjoints, so the same argument applied to
    $E(X)^\dagger$ gives
    $E(E(X)E(X)^\dagger)=E(X)E(X)^\dagger$.
    The vanishing of both Kadison--Schwarz gaps, again by the argument of
    Theorem~\ref{thm:fixedPointsStarSubalgebra}, places $E(X)$ in the
    multiplicative domain.
\end{proof}

\begin{theorem}[Choi--Effros projected product]
    \label{thm:choi_effros_sandwich}
    \lean{Kraus.choi_effros_sandwich}\leanok
    \uses{def:kraus_fixed_points, def:kraus_map}
    Under the same hypotheses, for all $X,Y\in\MN{D}$,
    $E(E(X)E(Y))=E(X)E(Y)$.
\end{theorem}

\begin{proof}\leanok
    \uses{thm:map_mem_multiplicativeDomain_of_idempotent}
    Theorem~\ref{thm:map_mem_multiplicativeDomain_of_idempotent} places
    $E(X)$ and $E(Y)$ in the multiplicative domain. Therefore
    $E(E(X)E(Y))=E(E(X))E(E(Y))=E(X)E(Y)$ by idempotence.
\end{proof}

The left, right, and symmetric absorption variants are
Theorems~\ref{thm:choi_effros_left}, \ref{thm:choi_effros_right}, and
\ref{thm:choi_effros_left_eq_right}, respectively.

\begin{definition}[Adjoint fixed points]\label{def:adjoint_fixed_points}
    \lean{Kraus.adjointFixedPoints}\leanok
    \uses{def:kraus_adjoint_map}
    The \emph{adjoint fixed-point set} is
    $\Fix(E^*)=\{X\in\MN{D}:E^*(X)=X\}$, where
    $E^*(X)=\sum_iK_i^\dagger XK_i$.
\end{definition}

\begin{theorem}[Fixed points form a $*$-subalgebra in the Heisenberg picture]
    \label{thm:adjointFixedPointsStarSubalgebra}
    \lean{Kraus.adjointFixedPointsStarSubalgebra, Kraus.fixedPoints_starSubalgebra}\leanok
    \uses{def:adjoint_fixed_points, def:tp_kraus}
    If the Kraus map is trace-preserving and has a positive definite
    fixed point $\rho>0$ with $E(\rho)=\rho$, then $\Fix(E^*)$ is a
    $*$-subalgebra of $\MN{D}$.
\end{theorem}

\begin{proof}\leanok
    \uses{thm:fixedPointsStarSubalgebra}
    Apply Theorem~\ref{thm:fixedPointsStarSubalgebra} to the adjoint
    family $\{K_i^\dagger\}$.
\end{proof}

\begin{definition}[Kraus commutant]\label{def:kraus_commutant}
    \lean{Kraus.krausCommutant, Kraus.krausCommutantStarSubalgebra}\leanok
    \uses{def:kraus_map}
    The \emph{Kraus commutant} is
    \begin{align}
        \{K_i,K_i^\dagger\}'
         & =\{X\in\MN{D}:[X,K_i]=[X,K_i^\dagger]=0\ \forall i\}.
        \label{eq:channel_kraus_commutant}
    \end{align}
    It is a $*$-subalgebra of $\MN{D}$.
\end{definition}

\begin{theorem}[Adjoint fixed points equal the Kraus commutant]
    \label{thm:adjointFixedPoints_eq_krausCommutant}
    \lean{Kraus.adjointFixedPointsStarSubalgebra_eq_krausCommutantStarSubalgebra}\leanok
    \uses{def:adjoint_fixed_points, def:kraus_commutant, def:tp_kraus}
    Under the hypotheses of Theorem~\ref{thm:adjointFixedPointsStarSubalgebra},
    the adjoint fixed-point $*$-subalgebra equals the Kraus commutant:
    \begin{align}
        \Fix(E^*)
         & =\{K_i,K_i^\dagger\}'.
        \label{eq:channel_adjoint_fix_commutant}
    \end{align}
    This is \cite[Theorem~6.13]{Wolf2012Quantum}.
\end{theorem}

\begin{proof}\leanok
    \uses{thm:adjointFixedPointsStarSubalgebra}
    For the inclusion ``$\supseteq$'', if $X$ commutes with all $K_i$ and
    $K_i^\dagger$, then
    $E^*(X)=\sum_iK_i^\dagger XK_i
        =X\sum_iK_i^\dagger K_i=X$ by trace preservation.
    For the inclusion ``$\subseteq$'', any $*$-subalgebra inside
    $\Fix(E^*)$ lies in the Kraus commutant, because for
    $X\in\Fix(E^*)$ with $X^\dagger X\in\Fix(E^*)$ the
    Kadison--Schwarz gap vanishes. This forces
    $XK_i=K_iX$ and $XK_i^\dagger=K_i^\dagger X$.
\end{proof}

\begin{theorem}[Kraus commutant inside the adjoint fixed points]
    \label{thm:kraus_commutant_greatest}
    \lean{Kraus.krausCommutantStarSubalgebra_isGreatest_adjointFixedPointStarSubalgebras}
    \leanok
    \uses{thm:adjointFixedPoints_eq_krausCommutant}
    Under the hypotheses above, the Kraus commutant is the largest
    $*$-subalgebra contained in the adjoint fixed-point set.
\end{theorem}

\begin{proof}\leanok
    By Theorem~\ref{thm:adjointFixedPoints_eq_krausCommutant}, every
    $*$-subalgebra inside the adjoint fixed points is contained in the Kraus
    commutant, while the Kraus commutant itself is such a $*$-subalgebra.
\end{proof}
